% !TEX program = lualatex
% ---------------------------------------------------------------------------------------------
%  Readback note 37: At most one primitive label for a Hecke eigensystem, by attached Galois representations
%
%  Build:  python3 tools/rb.py note <slug>
%
%  One unit of a Lean artifact, rewritten as ordinary mathematics. The prose is hand-written. Lean names,
%  line ranges and counts come from data.tex; quoted Lean statements are copied from the artifact by the build.
%
%  Rules of the text: standard terminology; every symbol is explained where it first appears and listed again
%  in the notation at the end; no letter has two meanings; a proof is complete or says what it takes for granted;
%  Lean names appear only in the LEAN lines and in the concordance.
%
%  The signature rule (the build enforces it): a function, relation or predicate that the artifact defines is
%  written only through its notation macro, which the build makes from the spec (`notation`) with one slot for
%  each explicit argument of the Lean definition. No argument is fixed and dropped, also not "for brevity" and
%  not for one section. Every argument and every hypothesis of a stated declaration stands inside its
%  statement, or in a Setting that the statement cites by number; only implicit and instance arguments (the
%  ambient objects) may stand in the conventions.
% ---------------------------------------------------------------------------------------------
\documentclass[11pt,a4paper]{article}
\ifdefined\pdfvariable\pdfvariable suppressoptionalinfo \numexpr 1+2+4+8\relax\fi
\usepackage[a4paper,hmargin=35mm,top=22mm,bottom=27mm,footskip=12mm,nohead]{geometry}

\usepackage{fontspec}
\directlua{
  luaotfload.add_fallback("gfallback", {"STIXTwoMath-Regular.otf:mode=node;", "DejaVuSans.ttf:mode=node;"})
  luaotfload.add_fallback("gmonofallback", {"DejaVuSansMono.ttf:mode=node;", "STIXTwoMath-Regular.otf:mode=node;"})
}
\setmainfont{NewCM10}[Extension=.otf,UprightFont=*-Book,ItalicFont=*-BookItalic,
  BoldFont=*-Bold,BoldItalicFont=*-BoldItalic,SmallCapsFont=*-Book,SmallCapsFeatures={Letters=SmallCaps},
  RawFeature={fallback=gfallback}]
\setsansfont{NewCMSans10}[Extension=.otf,UprightFont=*-Book,ItalicFont=*-BookOblique,
  BoldFont=*-Bold,BoldItalicFont=*-BoldOblique,RawFeature={fallback=gfallback}]
\setmonofont{JetBrainsMono}[Extension=.otf,UprightFont=*-Regular,ItalicFont=*-Italic,
  BoldFont=*-SemiBold,BoldItalicFont=*-SemiBoldItalic,Scale=MatchLowercase,
  RawFeature={-calt;-liga;fallback=gmonofallback}]
\usepackage{amsmath}
\usepackage{amsthm}
\usepackage[math-style=TeX]{unicode-math}
\setmathfont{NewCMMath-Book.otf}
\emergencystretch=2.5em\relax

\usepackage{xcolor}
\usepackage{tikz}
\usetikzlibrary{calc,arrows.meta,positioning,backgrounds}
\usepackage[protrusion=true,expansion=true]{microtype}
\usepackage{enumitem}
\usepackage{needspace}
\usepackage{booktabs}
\usepackage{array}
\usepackage{longtable}
\usepackage{fancyhdr}
\usepackage{lastpage}
\usepackage{multicol}
\usepackage[hidelinks,pdfusetitle]{hyperref}

\definecolor{ink}{HTML}{0F172A}
\definecolor{inksoft}{HTML}{475569}
\definecolor{inkfaint}{HTML}{94A3B8}
\definecolor{rulegrey}{HTML}{CBD5E1}
\definecolor{arith}{HTML}{C2410C}
\definecolor{est}{HTML}{7C3AED}
\definecolor{ana}{HTML}{0891B2}
\definecolor{ok}{HTML}{047857}
\definecolor{leancol}{HTML}{3730A3}

\title{At most one primitive label for a Hecke eigensystem, by attached Galois representations}
\author{grh-compact}
\date{}

% ------------------------------------------------------------------------------------- macros
\newcommand{\leanfont}{\ttfamily\fontsize{8.6}{10.5}\selectfont\color{leancol}}
\newcommand{\lean}[1]{\csname LN@#1\endcsname}
\newcommand{\fillin}[1]{{\color{arith}\itshape[#1]}}  % a place of the template that is not written yet; the build reports every one that is left
\newcommand{\leanquote}[1]{\csname LQ@#1\endcsname}
% generated by rb.py note; do not edit
\newcommand{\ConcDecls}{10}
\newcommand{\Pin}{72fead10}
\newcommand{\PinText}{that of commit {\leanfont 72fead10}}
\newcommand{\ArtifactName}{grh-compact}
\expandafter\newcommand\csname LN@isEigenformofhasHeckePolyAt\endcsname{\href{https://github.com/lsernau/grh-lean/blob/72fead102f663d5c7a407964a13dc0cf9d977207/grh-compact/src/New/GaloisRep/LabelUniqueness.lean\#L13-L16}{{\leanfont Data.\allowbreak{}isEigenform\_of\_hasHeckePolyAt}}}
\expandafter\newcommand\csname LQ@isEigenformofhasHeckePolyAt\endcsname{\leanquoteline{0}{lemma Data.isEigenform\_of\_hasHeckePolyAt \{M\} \{f : Gf \_ n → M → ℂ\} (h0 : f ≠ 0) \{Q\} (h : ∀ w}\leanquoteline{4}{(h : D.IsSplitOffLevel w), HasHeckePolyAt f w h.1 (Q w)) : IsEigenform f Q}}
\expandafter\newcommand\csname LN@occurringLabelssubsingleton\endcsname{\href{https://github.com/lsernau/grh-lean/blob/72fead102f663d5c7a407964a13dc0cf9d977207/grh-compact/src/New/GaloisRep/LabelUniqueness.lean\#L23-L59}{{\leanfont Data.\allowbreak{}occurringLabels\_subsingleton}}}
\expandafter\newcommand\csname LQ@occurringLabelssubsingleton\endcsname{\leanquoteline{0}{lemma Data.occurringLabels\_subsingleton \{v : \_ ↪ n\} \{S J χ' κ ν\} \{L : \_ →ₗ[ℂ] LevelPlaces.LevelForm}\leanquoteline{2}{S J χ' κ ν\} [Nonempty n] (hL : (⇑L).Injective) \{p\} \{F : \_ → \_\} (hF : ∀ y ∈ D.eigenPart V p, ∀ w}\leanquoteline{2}{(h : D.IsSplitOffLevel w), IsSphericalWith (rightRegular ..) h.1 (F w) (L y).toSmoothForm.1) x \{m}\leanquoteline{2}{b\} χ ψ (hk : ∀ y, HasParahoricType (s := m) (L y).toSmoothForm.1 (D.isConjApart\_chosen.2 x x) v b}\leanquoteline{2}{ψ χ fun g ↦ (L (D.principalRep V (Pi.mulSingle x g) y)).1.1) (hχ : χ.Injective) (hψ : ψ \textasciicircum{} P.Mω D.s}\leanquoteline{2}{= 1) (hM : ∀ l, χ l \textasciicircum{} P.Mω D.s = 1) : (D.occurringLabels V p x).Subsingleton}}
\expandafter\newcommand\csname LN@sphericalCompatstd\endcsname{\href{https://github.com/lsernau/grh-lean/blob/72fead102f663d5c7a407964a13dc0cf9d977207/grh-compact/src/Lib/GaloisRep/Main/SphericalAllRanks.lean\#L16-L88}{{\leanfont sphericalCompat\_std}}}
\expandafter\newcommand\csname LQ@sphericalCompatstd\endcsname{\leanquoteline{0}{lemma sphericalCompat\_std \{ℓ E n\} [ThePrime ℓ] [PadicIso ℓ] [Field E] [EmbeddedNumberField E]}\leanquoteline{2}{[CMType E]}\leanquoteline{2}{[Fintype n] [Nonempty n] [WeilData E] w : SphericalCompat E n (card n) w}}
\expandafter\newcommand\csname LN@isLabelledAt\endcsname{\href{https://github.com/lsernau/grh-lean/blob/72fead102f663d5c7a407964a13dc0cf9d977207/grh-compact/src/Lib/GaloisRep/LocalGlobal/LabelCompatibility.lean\#L24-L84}{{\leanfont Attached.\allowbreak{}isLabelledAt}}}
\expandafter\newcommand\csname LQ@isLabelledAt\endcsname{\leanquoteline{0}{lemma Attached.isLabelledAt \{ℓ\} [ThePrime ℓ] [PadicIso ℓ] \{E n\} [Nonempty n] [Field E] [CMType E]}\leanquoteline{2}{[EmbeddedNumberField E] [Fintype n] \{κ f P\} \{r :}\leanquoteline{2}{GalRep ℓ E n\} \{w A hw K\} [Field K] [Algebra (ResField w) K] \{T\} [IsFaithfulMulChar K T] \{s}\leanquoteline{2}{d c\} \{χ : Fin s → MulChar (ResField w) ℂ\} \{t\} [WeilData E] (hf : f ∈ (stdSpace E n).forms κ) (he :}\leanquoteline{2}{IsEigenform f P) (hr : Attached P r) (h : HasProductLabel f hw T d c χ t) (hA : IsPlaceOver A w) :}\leanquoteline{2}{IsFiniteTypeAt r A ∧ HasJumpsAt r A (jumps κ (Fintype.card n)) ∧}\leanquoteline{2}{IsLabelledAt A T d.factorial (fun s Q ↦ PolyKillsAt r A s Q ∧ (A.valuation}\leanquoteline{2}{ℓ = 1 → (r s).1.det = Q.roots.prod)) (c * ∏ l, (X - C (charName T (χ l)).1) \textasciicircum{} t l)}}
\expandafter\newcommand\csname LN@hasProductLabel\endcsname{\href{https://github.com/lsernau/grh-lean/blob/72fead102f663d5c7a407964a13dc0cf9d977207/grh-compact/src/Lib/GaloisRep/LocalGlobal/ParahoricTypeLabel.lean\#L31-L56}{{\leanfont HasParahoricType.\allowbreak{}hasProductLabel}}}
\expandafter\newcommand\csname LN@existsisTameExponent\endcsname{\href{https://github.com/lsernau/grh-lean/blob/72fead102f663d5c7a407964a13dc0cf9d977207/grh-compact/src/Lib/GaloisRep/LocalGlobal/TameExponent.lean\#L60-L79}{{\leanfont exists\_isTameExponent}}}
\expandafter\newcommand\csname LN@existsheckeSpliteqsmul\endcsname{\href{https://github.com/lsernau/grh-lean/blob/72fead102f663d5c7a407964a13dc0cf9d977207/grh-compact/src/Lib/Weil/Unramified/SphericalHecke.lean\#L255-L258}{{\leanfont exists\_heckeSplit\_eq\_smul}}}
\expandafter\newcommand\csname LN@isSplitofqprime\endcsname{\href{https://github.com/lsernau/grh-lean/blob/72fead102f663d5c7a407964a13dc0cf9d977207/grh-compact/src/Lib/General/Algebra/InertResidue.lean\#L31-L32}{{\leanfont isSplit\_of\_q\_prime}}}
\expandafter\newcommand\csname LN@scaleRoots\endcsname{\href{https://github.com/lsernau/grh-lean/blob/72fead102f663d5c7a407964a13dc0cf9d977207/grh-compact/src/New/Trace/BadShare.lean\#L35-L41}{{\leanfont IsPrimitiveLabel.\allowbreak{}scaleRoots}}}
\expandafter\newcommand\csname LN@isPrimitiveRootprodroots\endcsname{\href{https://github.com/lsernau/grh-lean/blob/72fead102f663d5c7a407964a13dc0cf9d977207/grh-compact/src/New/Trace/BadShare.lean\#L45-L51}{{\leanfont IsPrimitiveLabel.\allowbreak{}isPrimitiveRoot\_prod\_roots}}}
\newcommand{\ConcordanceRows}{\crow{isEigenformofhasHeckePolyAt}{\href{https://github.com/lsernau/grh-lean/blob/72fead102f663d5c7a407964a13dc0cf9d977207/grh-compact/src/New/GaloisRep/LabelUniqueness.lean\#L13-L16}{{\leanfont Data.\allowbreak{}isEigenform\_of\_hasHeckePolyAt}}}{LabelUniqueness.lean}{13--16}
\crow{occurringLabelssubsingleton}{\href{https://github.com/lsernau/grh-lean/blob/72fead102f663d5c7a407964a13dc0cf9d977207/grh-compact/src/New/GaloisRep/LabelUniqueness.lean\#L23-L59}{{\leanfont Data.\allowbreak{}occurringLabels\_subsingleton}}}{LabelUniqueness.lean}{23--59}
\crow{sphericalCompatstd}{\href{https://github.com/lsernau/grh-lean/blob/72fead102f663d5c7a407964a13dc0cf9d977207/grh-compact/src/Lib/GaloisRep/Main/SphericalAllRanks.lean\#L16-L88}{{\leanfont sphericalCompat\_std}}}{SphericalAllRanks.lean}{16--88}
\crow{isLabelledAt}{\href{https://github.com/lsernau/grh-lean/blob/72fead102f663d5c7a407964a13dc0cf9d977207/grh-compact/src/Lib/GaloisRep/LocalGlobal/LabelCompatibility.lean\#L24-L84}{{\leanfont Attached.\allowbreak{}isLabelledAt}}}{LabelCompatibility.lean}{24--84}
\crow{hasProductLabel}{\href{https://github.com/lsernau/grh-lean/blob/72fead102f663d5c7a407964a13dc0cf9d977207/grh-compact/src/Lib/GaloisRep/LocalGlobal/ParahoricTypeLabel.lean\#L31-L56}{{\leanfont HasParahoricType.\allowbreak{}hasProductLabel}}}{ParahoricTypeLabel.lean}{31--56}
\crow{existsisTameExponent}{\href{https://github.com/lsernau/grh-lean/blob/72fead102f663d5c7a407964a13dc0cf9d977207/grh-compact/src/Lib/GaloisRep/LocalGlobal/TameExponent.lean\#L60-L79}{{\leanfont exists\_isTameExponent}}}{TameExponent.lean}{60--79}
\crow{existsheckeSpliteqsmul}{\href{https://github.com/lsernau/grh-lean/blob/72fead102f663d5c7a407964a13dc0cf9d977207/grh-compact/src/Lib/Weil/Unramified/SphericalHecke.lean\#L255-L258}{{\leanfont exists\_heckeSplit\_eq\_smul}}}{SphericalHecke.lean}{255--258}
\crow{isSplitofqprime}{\href{https://github.com/lsernau/grh-lean/blob/72fead102f663d5c7a407964a13dc0cf9d977207/grh-compact/src/Lib/General/Algebra/InertResidue.lean\#L31-L32}{{\leanfont isSplit\_of\_q\_prime}}}{InertResidue.lean}{31--32}
\crow{scaleRoots}{\href{https://github.com/lsernau/grh-lean/blob/72fead102f663d5c7a407964a13dc0cf9d977207/grh-compact/src/New/Trace/BadShare.lean\#L35-L41}{{\leanfont IsPrimitiveLabel.\allowbreak{}scaleRoots}}}{BadShare.lean}{35--41}
\crow{isPrimitiveRootprodroots}{\href{https://github.com/lsernau/grh-lean/blob/72fead102f663d5c7a407964a13dc0cf9d977207/grh-compact/src/New/Trace/BadShare.lean\#L45-L51}{{\leanfont IsPrimitiveLabel.\allowbreak{}isPrimitiveRoot\_prod\_roots}}}{BadShare.lean}{45--51}}
% the notations of the defined functions: one slot for each explicit argument of the Lean declaration
\newcommand{\UnitMembers}{2}
\newcommand{\UnitInNote}{2}
\newcommand{\UnitProofs}{2}
\newcommand{\UnitSteps}{14}
\newcommand{\UnitAllSteps}{17}


\linespread{1.08}                       % 11 pt on 14.7 pt
\setlength{\parindent}{0pt}
\setlength{\parskip}{5.5pt plus 1.5pt minus 0.5pt}
\setlength{\abovedisplayskip}{7.5pt plus 2pt minus 3pt}
\setlength{\belowdisplayskip}{7.5pt plus 2pt minus 3pt}
\setlength{\abovedisplayshortskip}{3pt plus 2pt}
\setlength{\belowdisplayshortskip}{6pt plus 2pt minus 2pt}
\emergencystretch=1.5em
\clubpenalty=10000 \widowpenalty=10000 \displaywidowpenalty=10000
\relpenalty=7000 \binoppenalty=9000
\raggedbottom
\pagestyle{fancy}\fancyhf{}\renewcommand{\headrulewidth}{0pt}
\fancyfoot[C]{\sffamily\fontsize{7.6}{9}\selectfont\color{inkfaint}\thepage\,/\,\pageref*{LastPage}}
\setlist{topsep=2pt,partopsep=0pt,parsep=0pt,itemsep=4pt,leftmargin=2.1em,labelsep=0.55em}

\newcommand{\Q}{\mathbb{Q}}
\newcommand{\Z}{\mathbb{Z}}
\newcommand{\F}{\mathbb{F}}
\newcommand{\len}{\operatorname{len}}
\newcommand{\cok}{\operatorname{cok}}
\newcommand{\End}{\operatorname{End}}
\newcommand{\Hom}{\operatorname{Hom}}
\newcommand{\Gal}{\operatorname{Gal}}
\newcommand{\diag}{\operatorname{diag}}
\newcommand{\kw}[1]{\emph{#1}}
\newcommand{\m}[1]{\mbox{$#1$}}        % a short formula that must not be broken across lines
\newcommand{\head}[1]{\par\addvspace{17pt plus 4pt minus 2pt}\Needspace{6\baselineskip}%
  \noindent{\fontsize{13.2}{16}\selectfont\bfseries\color{ink}#1}\par\nopagebreak\vspace{-1pt}}
\newcommand{\sect}[2]{\refstepcounter{section}\label{#2}\par\addvspace{19pt plus 4pt minus 2pt}\Needspace{7\baselineskip}%
  \noindent\llap{{\fontsize{13.2}{16}\selectfont\bfseries\color{inkfaint}\thesection}\hspace{9pt}}%
  {\fontsize{13.2}{16}\selectfont\bfseries\color{ink}#1}\par\nopagebreak\vspace{-1pt}}
\newcommand{\subhead}[1]{\par\addvspace{9pt plus 2pt}\Needspace{4\baselineskip}\noindent\textbf{#1}\enspace\ignorespaces}
% where a statement lives in the Lean archive
\newcommand{\leantag}{{\sffamily\bfseries\fontsize{6.3}{8}\selectfont\color{leancol}\addfontfeatures{LetterSpace=7}LEAN}}
\newcommand{\sep}{\unskip\nobreak\hspace{5pt}{\color{inkfaint}\textperiodcentered}\hspace{5pt plus 2pt}\ignorespaces}
\newcommand{\leanline}[1]{\par\nopagebreak\vspace{-2.5pt}{\fontsize{8.9}{12.4}\selectfont\color{inksoft}\raggedright
  \noindent\llap{\leantag\hspace{9pt}}#1\par}\vspace{1pt}}
% a displayed statement in a tinted box: #1 colour, #2 text
\newcommand{\statement}[2]{\par\addvspace{7pt}\begingroup\setlength{\fboxsep}{8.5pt}\setlength{\fboxrule}{0.6pt}%
  \noindent\fcolorbox{#1!30!white}{#1!3!white}{\parbox{\dimexpr\linewidth-2\fboxsep-2\fboxrule\relax}{#2}}%
  \endgroup\par\addvspace{7pt}}
\newcommand{\nt}[2]{\par\noindent\hangindent=1.1em\hangafter=1 #1\enspace{\color{inksoft}#2}\par\vspace{1pt}}
\newcommand{\legdot}[1]{\tikz[baseline=-0.62ex]\fill[#1] (0,0) circle (2.5pt);\hspace{3pt}}
\newcommand{\crow}[4]{\csname CS@#1\endcsname & #2 & \textcolor{inksoft}{#3} & \textcolor{inksoft}{#4}\\}
\newcommand{\leanquoteline}[2]{\par\noindent\hangindent=\dimexpr#1\fontdimen2\font+2.2em\relax\hangafter=1
  \hspace*{#1\fontdimen2\font}#2}
\newcommand{\quoted}[1]{\par\vspace{3pt}{\ttfamily\fontsize{7.5}{10.4}\selectfont\color{leancol}\raggedright
  \setlength{\parskip}{0pt}\leanquote{#1}\par}\vspace{4pt}}
\newcommand{\bx}[3]{{\sffamily\bfseries\fontsize{7.3}{9.2}\selectfont\color{#1!82!black}#2}\\[1.5pt]#3}
\newcommand{\cs}[2]{\expandafter\newcommand\csname CS@#1\endcsname{#2}}

% theorems
\newtheoremstyle{plainx}{7pt plus 2pt}{7pt plus 2pt}{\itshape}{}{\bfseries}{.}{0.5em}{}
\newtheoremstyle{uprightx}{7pt plus 2pt}{7pt plus 2pt}{}{}{\bfseries}{.}{0.5em}{}
\theoremstyle{plainx}
\newtheorem{theorem}{Theorem}[section]
\newtheorem{lemma}[theorem]{Lemma}
\newtheorem{proposition}[theorem]{Proposition}
\newtheorem{corollary}[theorem]{Corollary}
\theoremstyle{uprightx}
\newtheorem{definition}[theorem]{Definition}
\newtheorem{setting}[theorem]{Setting}          % hypotheses shared by several statements; each of them cites it by number
\newtheorem{remark}[theorem]{Remark}
\newtheorem{example}[theorem]{Example}
\renewcommand{\qedsymbol}{{\setlength{\fboxsep}{0pt}\setlength{\fboxrule}{0.5pt}\fbox{\rule{0pt}{0.46em}\hspace{0.46em}}}}
% what holds throughout the note: the ambient objects only (implicit and instance arguments), never a hypothesis
\newenvironment{conventions}{\par\addvspace{5pt}\noindent\textbf{Conventions.}\enspace\ignorespaces}{\par\addvspace{3pt}}
\makeatletter
\renewenvironment{proof}[1][Proof]{\par\pushQED{\qed}\normalfont\topsep0pt\partopsep0pt
  \trivlist\item[\hskip\labelsep\itshape #1.]\ignorespaces}{\popQED\endtrivlist\@endpefalse}
\makeatother

% where each declaration of the concordance is treated: \cs{key}{place}, the place naming a \label or a lettered item
\cs{occurringLabelssubsingleton}{Theorem 1}
\cs{isEigenformofhasHeckePolyAt}{Lemma~\ref{l:eigen}}
\cs{sphericalCompatstd}{Input 1}
\cs{isLabelledAt}{Input 2}
\cs{hasProductLabel}{Input 3}
\cs{existsisTameExponent}{Input 4}
\cs{existsheckeSpliteqsmul}{Input 4}
\cs{isSplitofqprime}{Input 4}
\cs{scaleRoots}{Input 5}
\cs{isPrimitiveRootprodroots}{Input 5}
\newcommand{\cV}{{\mathcal V}}
\newcommand{\fs}{{\mathfrak s}}
\newcommand{\GL}{\mathrm{GL}}

\begin{document}
\color{ink}

{\raggedright
{\sffamily\bfseries\fontsize{7.2}{9}\selectfont\color{leancol}\addfontfeatures{LetterSpace=8}READBACK NOTE 37\par}
\vspace{5pt}
{\fontsize{21.5}{25.5}\selectfont\bfseries At most one primitive label for a Hecke eigensystem, by attached Galois representations\par}
\vspace{5pt}
{\fontsize{12.4}{16}\selectfont\itshape\color{inksoft} One unit of the Lean artifact \ArtifactName, read back into
mathematics\par}}
\vspace{8pt}
{\color{rulegrey}\hrule height 0.6pt}
\vspace{7pt}
\thispagestyle{empty}

\textbf{Summary.} At a chosen prime the forms of principal level carry a representation of a finite general
linear group, and a primitive label names a cuspidal type of that group. This note proves that a Hecke
eigensystem can meet at most one primitive label at a chosen prime, as soon as the forms can be lifted
injectively to a definite unitary group in a way that respects the Hecke operators away from the level and the
action of a parabolic subgroup at the prime. The proof passes to Galois representations: a lifted form is an
eigenform, so a $2$-adic Galois representation is attached to its Hecke polynomials; local-global compatibility at
tame level reads the label off the action of tame inertia; and two different primitive labels would force the
eigenvalues of one inertia element to be roots of unity of small order, against the determinant.

\statement{leancol}{\fontsize{9.6}{13.2}\selectfont\textbf{What this note is.} The artifact \ArtifactName{} is a
body of Lean code: the compact edition of a formal proof, due to L.~Alpöge with the AI system Claude, of the
Riemann hypothesis and of its generalisation to Dirichlet $L$-functions. That headline is a claim. The authors
report that every file compiles and that Lean's kernel accepted the proofs; an independent replay is recorded in
the handout's record of checks when it reports. No mathematician has read the argument through, and nothing here
says that the Riemann hypothesis is proved. This note treats the unit \emph{At most one primitive label per
eigensystem at a chosen place}: all \UnitInNote{} of its \UnitMembers{} declarations, a lemma and a theorem. Both
are proved in full from the inputs. Taken for granted, and stated as Inputs~1 to~5, are: the existence of Galois
representations attached to eigenforms on definite unitary groups, and local-global compatibility at tame level,
both in the weak forms in which the artifact's background library proves them; three further facts of that
library; and facts about the data and about primitive labels from other units. Every statement was compared with
its Lean text, argument by argument. The Lean text read is \PinText; Appendix~A is the concordance of
\ConcDecls{} declarations.}

\sect{The statement}{s:intro}
The setting and the notions are those of the note \emph{Good forms and rows}, recalled in
Setting~\ref{s:data}; algebraic forms, Hecke polynomials and block triangular matrices are explained in
§\ref{s:setting}.

\statement{arith}{\phantomsection\label{t:1}\textbf{Theorem 1} (at most one primitive label).\enspace
\emph{In Setting~\ref{s:data}, let $\fs$ be a family of signs on $I(D)$ and $\cV$ a space over $E(D)$ with index
set $I(D)$ and signs~$\fs$. Let $I'$ be a non-empty finite set, \m{\nu\colon I(D)\to I'} an injection, and $L$ a
linear map from \m{\mathrm{Prin}(D,\cV)} to a space of algebraic modular forms on the standard space with index
set~$I'$, in the library's sense: such a space is given by a weight in two parts, a system of places of $E(D)$, a
subgroup of the product of the groups \m{\GL_{I'}(\F_w)} over these places, and a character of that subgroup. Let
$p$ be an eigensystem, \m{(\Upsilon_w)_w} a family of functions \m{\mathbb C\to\mathbb C} indexed by the primes
of~$E(D)$, $v$ a chosen prime, $n_1$ a natural number, \m{\beta\colon I'\to\{0,\dots,n_1\}} a map,
\m{\chi=(\chi_1,\dots,\chi_{n_1})} a family of multiplicative characters of~$\F_v$, and $\psi$ a multiplicative
character of~$\F_v$. Suppose:}
\begin{enumerate}[label=\textup{(\roman*)},leftmargin=2.3em,itemsep=1pt]
\item \emph{$L$ is injective;}
\item \emph{for every \m{y\in\mathrm{Prin}(D,\cV)[p]} and every split prime $w$ of $E(D)$ such that neither $w$ nor
$\bar w$ is chosen, the function $L(y)$ is spherical at $w$ with~$\Upsilon_w$: it is fixed by right translation
by the elements of $K_0(I')$ with component~$1$ at the primes other than $w$ and~$\bar w$, and for every
\m{X\in\mathbb C}
\[\sum_{n=0}^{\#I'}(-X)^n\,(\#\F_w)^{n(n-1)/2}\;\mathbb T^{w,n}L(y)=\Upsilon_w(X)\,L(y);\]}
\item \emph{for every \m{y\in\mathrm{Prin}(D,\cV)}, the function $L(y)$ has parahoric type at~$v$ for
\m{(\nu,\beta,\psi,\chi)} and the functions \m{x\mapsto L\bigl(\mathrm{rep}_{D,\cV}(x_{(v)})\,y\bigr)}: that is,
\m{\beta^{-1}(0)=\nu(I(D))}, and for every \m{x\in\GL_{I(D)}(\F_v)} and every \m{k\in\GL_{I'}(O_v)} whose reduction
$\bar k$ is block upper triangular for $\beta$ with block~$0$ equal to~$x$,
\[R(k^{(v)})\,L(y)=\prod_{n=1}^{n_1}\chi_n(\det\bar k_{nn})\cdot\psi(\det x)\cdot L\bigl(\mathrm{rep}_{D,\cV}(x_{(v)})\,y\bigr);\]}
\item \emph{the characters \m{\chi_1,\dots,\chi_{n_1}} are pairwise distinct;}
\item \emph{\m{\psi^{M_\omega(D)}=1};}
\item \emph{\m{\chi_n^{M_\omega(D)}=1} for every~$n$.}
\end{enumerate}
\emph{Then the set \m{\mathrm{Occ}(D,\cV,p,v)} has at most one element.}}
\leanline{\lean{occurringLabelssubsingleton}}

The theorem says nothing about existence: if \m{\mathrm{Prin}(D,\cV)[p]} is zero, no label occurs. It is applied
twice in the note \emph{Good forms and rows}: to the space of the rows, with $L$ the chain of $m(D)$ theta lifts
to a definite space, and to the source of each single lift, with $L$ that lift. In both cases the lift carries
the cuspidal type at~$v$ into a type induced from a parabolic subgroup, which is what (iii) expresses.

\sect{The setting}{s:setting}
\begin{conventions}\phantomsection\label{c:conventions}
The natural numbers are \m{0,1,2,\dots}. The parameters \m{P=(N_0,d)} of the proof are fixed. For a number field
$E$: a prime of~$E$ is a non-zero prime ideal of its ring of integers; $\F_w$, $E_w$, $O_w$ are the residue
field, the completion and its ring of integers at a prime~$w$; for a CM field, \m{z\mapsto\bar z} is complex
conjugation, and a prime $w$ is split if \m{\bar w\ne w}. For a finite set~$I'$, \m{\GL_{I'}} of a ring is the group
of invertible matrices over it indexed by~$I'$; \m{G_f(I')} is the group of the matrices $g$ over the finite
adeles of~$E$, indexed by~$I'$, with \m{g^*g=1}, where $g^*$ is the conjugate transpose; \m{K_0(I')} is its
subgroup of the matrices with integral entries; for a split prime $w$ and \m{k\in\GL_{I'}(E_w)}, \m{k^{(w)}\in G_f(I')}
has component $k$ at~$w$ and $1$ at the primes other than $w$ and~$\bar w$; and $R$ is right translation of
functions on~$G_f(I')$, \m{(R(g)f)(g')=f(g'g)}. $\bar{\mathbb Q}$ is an algebraic closure of~$\mathbb Q$ that
contains~$E$. Labels, primitive labels, the polynomial \m{c^{[u]}} whose roots are $u$ times those of~$c$, the
injective character $T$ of the multiplicative group of an algebraic closure $\bar{\F}$ of a finite field~$\F$,
the unit $x_\chi$ that names a multiplicative character~$\chi$, Green's class functions $\Xi_c$ and their
induction product~$*$, and the idempotents \m{e_{I,c}} are those of the note \emph{Primitive labels}. Letters of
the handout's overview that are not displayed with the data~$D$ as an argument (among them $t$, $\lambda$,
$\Upsilon$, $\delta$, $\zeta$ and~$j$) have, in this note, only the meaning given in its notation list.
\end{conventions}

\begin{setting}[the data and its forms]\label{s:data}
$D$ is admissible data for the parameters~$P$ (Chapter~III of the handout). The following objects are those of
the setting and of the definitions of the note \emph{Good forms and rows}: the CM field $E(D)$, which
contains a primitive root of unity of order~$5040$; the chosen primes, none of which is conjugate to a chosen
prime; the rank \m{r(D)>1}; the row index set $I(D)$ with $r(D)$ elements and the group
\m{\mathsf G(D)=\prod_v\GL_{I(D)}(\F_v)} over the chosen primes; the integer \m{M_\omega(D)}, for which
\m{M_\omega(D)^2} divides \m{\#\F_v-1} and \m{0<M_\omega(D)<\#\F_v-1} at every chosen prime~$v$; the fact that
$\#\F_v$ is an odd prime for chosen~$v$; for signs $\fs$ on $I(D)$ and a space $\cV$ over $E(D)$ with index set
$I(D)$ and signs~$\fs$, the space \m{\mathrm{Prin}(D,\cV)} of forms of principal level at the chosen primes and the
representation \m{\mathrm{rep}_{D,\cV}} of \m{\mathsf G(D)} on it; the Hecke operators \m{\mathbb T^{w,n}} at the
split primes $w$ with neither $w$ nor $\bar w$ chosen; eigensystems $p$ and the eigenparts
\m{\mathrm{Prin}(D,\cV)[p]}; and, for a chosen prime~$v$, the set \m{\mathrm{Occ}(D,\cV,p,v)} of the primitive
labels \m{c\in\F_v[X]} of degree $r(D)$ for which \m{e_{I(D),c}}, acting through the factor $v$ of
\m{\mathsf G(D)}, does not kill \m{\mathrm{Prin}(D,\cV)[p]}. For \m{x\in\GL_{I(D)}(\F_v)}, \m{x_{(v)}\in\mathsf G(D)}
has component $x$ at $v$ and $1$ elsewhere.
\end{setting}

\subhead{Functions on $G_f(I')$: Hecke polynomials, eigenforms, types, labels.} These notions belong to the
background library. Let $E$ be a CM field, $I'$ a finite set, $\mathcal W$ a set, and $f$ a function on $G_f(I')$
with values in the complex functions on~$\mathcal W$.
\begin{enumerate}[label=\textup{(\alph*)},leftmargin=2.3em,itemsep=2pt]
\item For a split prime~$w$ and a natural number~$n$, \m{\mathbb T^{w,n}f} is the sum of the right translates of
$f$ by the \m{k^{(w)}}, $k$ running over representatives of the cosets modulo \m{\GL_{I'}(O_w)} in the double coset
of \m{\mathrm{diag}(\varpi,\dots,\varpi,1,\dots,1)}, with $n$ entries equal to a uniformizer~$\varpi$.
\item $f$ \emph{has Hecke polynomial $Q$ at a split prime~$w$}, for a complex polynomial~$Q$, if $f$ is fixed by
right translation by the elements of $K_0(I')$ with component $1$ at the primes other than $w$ and~$\bar w$, and
there are complex numbers \m{t_0,t_1,\dots} with \m{\mathbb T^{w,n}f=t_nf} for all~$n$ and
\m{Q=\sum_{n=0}^{\#I'}(-1)^n(\#\F_w)^{n(n-1)/2}\,t_n\,X^{\#I'-n}}.
\item $f$ \emph{is an eigenform for a family \m{(\mathrm{Hk}_w)_w} of complex polynomials}, indexed by the primes
of~$E$, if \m{f\ne0} and for all but finitely many primes~$w$: if the norm $\#\F_w$ is a prime number, then $w$ is
split and $f$ has Hecke polynomial $\mathrm{Hk}_w$ at~$w$.
\item \emph{Algebraic modular forms} on the standard space with index set $I'$, of a weight~$\kappa$: functions
$f$ as above, for a finite set $\mathcal W$ fixed by the weight, that are invariant under right translation by an
open subgroup and satisfy \m{f(\gamma g)=\mathrm{alg}_\kappa(\gamma)f(g)} for all rational points $\gamma$ of the
definite unitary group, with $\mathrm{alg}_\kappa$ the algebraic representation of the weight.
\item Let \m{\beta\colon I'\to\{0,\dots,n_1\}} be a map. A matrix $\bar k$ indexed by $I'$ is \emph{block upper
triangular for~$\beta$} if its entry at \m{(\iota,\iota')} vanishes whenever \m{\beta(\iota')<\beta(\iota)}; its
$n$-th diagonal block $\bar k_{nn}$ is its part with rows and columns in \m{\beta^{-1}(n)}. With an injection
$\nu$ of a set $I$ onto \m{\beta^{-1}(0)}, block~$0$ is read as a matrix indexed by~$I$.
\item Let $w$ be a split prime, \m{c\in\F_w[X]}, \m{\chi_1,\dots,\chi_{n_1}} multiplicative characters of $\F_w$ and
\m{t_1,\dots,t_{n_1}} natural numbers. Then $f$ \emph{has the product label \m{(c;\chi;t)} of degree $\delta$
at~$w$} if: $c$ is a label of degree~$\delta$; the polynomial \m{c\cdot\prod_n(X-x_{\chi_n})} is separable;
\m{\delta+\sum_nt_n=\#I'}; $f$ is fixed by \m{R(k^{(w)})} for the \m{k\in\GL_{I'}(O_w)} that reduce to~$1$; and
\m{\sum_{x}\overline{\Xi(x)}\,R(\tilde x^{(w)})f\ne0}, the sum over \m{x\in\GL_{I'}(\F_w)} with lifts
\m{\tilde x\in\GL_{I'}(O_w)}, where $\Xi$ is the induction product of $\Xi_c$ and of the functions
\m{\chi_n\circ\det} in size~$t_n$.
\end{enumerate}

\subhead{Galois representations.} Let $E$ be a CM field inside $\bar{\mathbb Q}$, $\ell$ a prime number and
\m{\iota\colon\mathbb C\to\bar{\mathbb Q}_\ell} an isomorphism of fields. A \emph{place of $\bar{\mathbb Q}$ over
a prime $w$} of $E$ is a valuation ring $\mathfrak P$ of $\bar{\mathbb Q}$ whose maximal ideal meets the integers
of $E$ in~$w$; its inertia group consists of the \m{\sigma\in\Gal(\bar{\mathbb Q}/E)} that preserve $\mathfrak P$ and
act trivially on its residue field, which is an algebraic closure $\bar{\F}_w$ of~$\F_w$.
\begin{enumerate}[label=\textup{(\alph*)},leftmargin=2.3em,itemsep=2pt,start=7]
\item \emph{Tame exponents.} Let $a$ and \m{\delta\ne0} be natural numbers and \m{q=(\#\F_w)^\delta}. An element
$\sigma$ of the inertia group of $\mathfrak P$ \emph{has tame exponent $a$ of level~$\delta$} if there are a
uniformizer $\varpi$ of $E$ at $w$ and \m{z\in\bar{\mathbb Q}} with \m{z^{q-1}=\varpi} such that the reduction
\m{\zeta\in\bar{\F}_w} of the root of unity \m{\sigma(z)/z} satisfies \m{T(\zeta)=\exp(2\pi i\,a/(q-1))}.
\item \emph{Attached representations.} A homomorphism \m{\rho\colon\Gal(\bar{\mathbb Q}/E)\to\GL_{I'}(\bar{\mathbb Q}_\ell)}
is \emph{attached to} a family \m{(\mathrm{Hk}_w)_w} of complex polynomials if it is continuous, if the trace
form \m{(a,b)\mapsto\mathrm{tr}(ab)} is non-degenerate on the linear span of its image (so $\rho$ is semisimple),
and if for every real \m{\eta>0} the following set of Frobenius data is \emph{large}: the triples
\m{(w,\mathfrak P,\sigma)} of a prime $w$ of prime norm, a place $\mathfrak P$ over it and a geometric Frobenius
element $\sigma$ at~$\mathfrak P$ such that every coefficient of the characteristic polynomial of
\m{\rho(\sigma)} differs from the corresponding coefficient of \m{\iota(\mathrm{Hk}_w)} by at most $\eta$ in the
$\ell$-adic absolute value. A set of Frobenius data is large if every non-empty open subset of the Galois group
contains a non-empty open subset $U$ such that, for all but finitely many~$w$, all data \m{(w,\mathfrak P,\sigma)}
with \m{\sigma\in U} belong to the set.
\end{enumerate}
So ``attached'' is an approximate notion: it does not say that the characteristic polynomial of a Frobenius
element is the Hecke polynomial, only that this holds up to every $\ell$-adic precision on a set of Frobenius
elements that is dense in a strong sense. By the theorem of Chebotarev (which the library proves in its existence
form) and continuity, this still determines the characteristic polynomial of $\rho$ on the whole group, and
hence $\rho$ up to isomorphism by the theorem of Brauer and Nesbitt; the proof below needs neither, because
Input~2 holds for every attached representation.

\sect{Inputs}{s:inputs}
\statement{inksoft}{\textbf{Inputs from the background library.} $E$ is a CM field inside $\bar{\mathbb Q}$ with
Weil data in the library's sense (in particular it contains a primitive root of unity of order~$5040$); $I'$ is a
non-empty finite set; $\ell$ is a prime and \m{\iota\colon\mathbb C\to\bar{\mathbb Q}_\ell} an isomorphism of fields.

\vspace{4pt}\phantomsection\label{i:1}\textbf{Input 1} (attached Galois representations exist).\enspace
\emph{Let \m{(\mathrm{Hk}_w)_w} be a family of complex polynomials and $f$ an algebraic modular form of some weight
on the standard space with index set~$I'$ that is an eigenform for \m{(\mathrm{Hk}_w)_w}. Then there is a
homomorphism \m{\rho\colon\Gal(\bar{\mathbb Q}/E)\to\GL_{I'}(\bar{\mathbb Q}_\ell)} attached to \m{(\mathrm{Hk}_w)_w}.}

\vspace{3pt}\phantomsection\label{i:2}\textbf{Input 2} (local-global compatibility at tame level).\enspace
\emph{Let $f$ be an algebraic modular form of some weight on the standard space with index set~$I'$ that is an
eigenform for \m{(\mathrm{Hk}_w)_w}, and let $\rho$ be \emph{any} homomorphism attached to \m{(\mathrm{Hk}_w)_w}.
Let $w$ be a split prime at which $f$ has the product label \m{(c;\chi;t)} of degree~$\delta$, $\mathfrak P$ a place
of $\bar{\mathbb Q}$ over~$w$, and $\sigma$ an element of the inertia group of $\mathfrak P$ that has tame
exponent $a$ of level~$\delta!$. Let
\m{Q=\prod_x\bigl(X-\iota(T(x)^a)\bigr)}, the product over the roots $x$ in $\bar{\F}_w$, with multiplicity, of
\m{c\cdot\prod_n(X-x_{\chi_n})^{t_n}}. If $\ell$ is not the residue characteristic of~$w$, then
\m{Q(\rho(\sigma))=0}, and \m{\det\rho(\sigma)} is the product of the roots of~$Q$.}

\vspace{3pt}\phantomsection\label{i:3}\textbf{Input 3} (type implies label).\enspace \emph{Let $w$ be a split
prime, $I$ a finite set, and let \m{\GL_I(\F_w)} act linearly on a complex vector space, written
\m{(x,y)\mapsto x\cdot y}; let $y$ be a vector of it, and $L$ an injective linear map from that space to the
functions on $G_f(I')$ with values in the functions on a set. Let \m{c\in\F_w[X]} with \m{e_{I,c}\cdot y\ne0}. Let $\nu$, $\beta$, $\psi$, \m{\chi_1,\dots,\chi_{n_1}} be
as in Theorem~1\textup{(iii)}, with the $\chi_n$ pairwise distinct, and suppose that $L(y)$ has parahoric type
at~$w$ for \m{(\nu,\beta,\psi,\chi)} and the functions \m{x\mapsto L(x\cdot y)}. If no $x_{\chi_n}$ is a root of
\m{c^{[x_\psi]}}, then $L(y)$ has the product label \m{(c^{[x_\psi]};\chi;t)} of degree~$\#I$ at~$w$, with
\m{t_n=\#\beta^{-1}(n)}.}

\vspace{3pt}\phantomsection\label{i:4}\textbf{Input 4} (tame inertia; Hecke polynomials).\enspace
\emph{\textup{(a)} For every prime $w$ of~$E$, every place $\mathfrak P$ of $\bar{\mathbb Q}$ over $w$ and every
natural number~$\delta$, some element of the inertia group of $\mathfrak P$ has tame exponent~$1$ of
level~$\delta!$. \textup{(b)} Let $w$ be a split prime and \m{\Upsilon\colon\mathbb C\to\mathbb C} a function. There
are complex numbers \m{t_0,t_1,\dots} such that \m{\mathbb T^{w,n}f=t_nf} for all~$n$ and every function $f$ on
$G_f(I')$ that is spherical at $w$ with~$\Upsilon$ in the sense of Theorem~1\textup{(ii)}. \textup{(c)} A prime
of~$E$ whose norm is a prime number is split.}}
\leanline{\lean{sphericalCompatstd}\sep \lean{isLabelledAt}\sep \lean{hasProductLabel}\sep
\lean{existsisTameExponent}\sep \lean{existsheckeSpliteqsmul}\sep \lean{isSplitofqprime}}

\statement{inksoft}{\phantomsection\label{i:5}\textbf{Input 5} (primitive labels; the two lemmas of the note
\emph{Primitive labels}).\enspace \emph{Let $\F$ be a finite field, \m{c\in\F[X]} a primitive label of degree~$r$.
\textup{(a)} If \m{u\in\F^\times} and $M$ is a natural number with \m{u^M=1} and \m{M^2} dividing \m{\#\F-1}, then
\m{c^{[u]}} is a primitive label of degree~$r$. \textup{(b)} The product of the roots of $c$ in $\bar{\F}$ is a
primitive root of unity of order \m{\#\F-1}.}}
\leanline{\lean{scaleRoots}\sep \lean{isPrimitiveRootprodroots}}

\subhead{On Inputs 1 and 2.} These are the two deep inputs. The references and their relation to the formal
statements are those of the catalogue of this handout and of its second reading, which compared the statements
with the sources.

\emph{Input 1} is a weaker consequence of the construction of Galois representations for automorphic
representations of unitary groups: G.~Chenevier and M.~Harris, \emph{Construction of automorphic Galois
representations,~II}, Camb. J. Math. 1 (2013), Theorem~3.2.3; T.~Barnet-Lamb, T.~Gee, D.~Geraghty and R.~Taylor,
\emph{Potential automorphy and change of weight}, Ann. of Math. 179 (2014), Theorem~2.1.1; with base change from
the unitary group by J.-P.~Labesse, \emph{Changement de base CM et séries discrètes}, Corollaire~5.3. It is weaker
in three ways: only the standard definite hermitian space over a CM field with Weil data is treated; ``attached''
has the approximate meaning of (h) above, in place of the equality of the characteristic polynomial of Frobenius
with the Hecke polynomial at almost all places; and the transport between eigenforms on the unitary group,
automorphic representations of the general linear group and their weights is part of the library's own argument.
The formal statement has a second half, on the local behaviour of attached representations at places where the
form is spherical, which is not used in this unit.

\emph{Input 2} goes in the direction of local-global compatibility, for which the library's file names
M.~Harris and R.~Taylor, \emph{The geometry and cohomology of some simple Shimura varieties} (2001), R.~Taylor and
T.~Yoshida, \emph{Compatibility of local and global Langlands correspondences}, J. Amer. Math. Soc. 20 (2007),
Theorem~A, and T.~Barnet-Lamb, T.~Gee, D.~Geraghty and R.~Taylor, \emph{Local-global compatibility for $l=p$,~II},
Ann. Sci. ENS 47 (2014), Theorem~1.1. Those theorems identify the Weil--Deligne representation of
$\rho$ at $w$ with the local Langlands parameter of the local component. Input~2 retains much less: only split
places, only the restriction to tame inertia, only that a polynomial kills $\rho(\sigma)$ together with the value
of the determinant, and with the local component described by a finite type (Green's character induced with
characters of the determinant) in place of a parameter. The second reading leaves open the identification of
this formal statement with a published one. According to the header of the library's file (I read the statement
and not its proof), it is reduced there to the spherical case by cyclic base change. In classical terms it says that tame inertia acts through the sum of the characters $\theta_x$, $x$ the
roots of the label, composed with the fundamental character of level~$\delta!$, as the depth-zero local Langlands
correspondence predicts. The formal statement also contains assertions on finite ramification and on
Hodge--Tate weights, and a clause for $\ell$ equal to the residue characteristic, none of which is used here.

Input~3 is Frobenius reciprocity for the parahoric subgroup together with the orthogonality of Green's
characters; Input~4(a) is Kummer theory of the tame quotient of inertia; Input~4(b) says that the values of the
Hecke polynomial determine the eigenvalues; Input~4(c) holds because a prime fixed by complex conjugation has
square norm when $E$ contains a primitive root of unity of order~$5040$. The library proves these four without a
named source.

\sect{Proofs}{s:proofs}
\begin{lemma}[eigenforms from Hecke polynomials off the level]\label{l:eigen}
In Setting~\ref{s:data}, let $I'$ be a finite set, $\mathcal W$ a set, $f$ a function on $G_f(I')$ with values in
the complex functions on~$\mathcal W$, and \m{(\mathrm{Hk}_w)_w} a family of complex polynomials indexed by the
primes of~$E(D)$. Suppose that \m{f\ne0}, and that for every split prime $w$ of $E(D)$ such that neither $w$ nor
$\bar w$ is chosen, $f$ has Hecke polynomial $\mathrm{Hk}_w$ at~$w$. Then $f$ is an eigenform for
\m{(\mathrm{Hk}_w)_w}.
\end{lemma}
\begin{proof}
There are finitely many chosen primes, so all but finitely many primes $w$ are such that neither $w$ nor
$\bar w$ is chosen. If such a $w$ has prime norm, it is split by Input~4(c), and the hypothesis applies to it.
\end{proof}
\leanline{\lean{isEigenformofhasHeckePolyAt}}

\begin{proof}[Proof of Theorem 1]
Let \m{c_1,c_2\in\mathrm{Occ}(D,\cV,p,v)}; we show \m{c_1=c_2}. Write \m{q=\#\F_v} and \m{r=r(D)}, and let $T$ be
the injective character of \m{\bar{\F}_v^\times}. Take \m{\ell=2}; an isomorphism of fields
\m{\iota\colon\mathbb C\to\bar{\mathbb Q}_2} exists, because both fields are algebraically closed, of
characteristic~$0$ and of the cardinality of the continuum. Fix a place $\mathfrak P$ of $\bar{\mathbb Q}$
over~$v$ and, by Input~4(a), an element $\sigma$ of its inertia group with tame exponent~$1$ of level~$r!$. As $q$
is odd, $2$ is not the residue characteristic of~$v$.

For \m{j\in\{1,2\}} choose \m{y_j\in\mathrm{Prin}(D,\cV)[p]} such that \m{e_{I(D),c_j}}, acting through the factor
$v$ of \m{\mathsf G(D)}, does not kill~$y_j$; put \m{f_j=L(y_j)} and \m{c_j'=c_j^{[x_\psi]}}.

\emph{Step 1: $c_j'$ is a primitive label of degree~$r$.} By (v), \m{x_\psi^{M_\omega(D)}=1}, and
\m{M_\omega(D)^2} divides \m{q-1}; apply Input~5(a).

\emph{Step 2: $f_j$ has the product label \m{(c_j';\chi;t)} of degree $r$ at~$v$, with \m{t_n=\#\beta^{-1}(n)}.}
This is Input~3 for the representation \m{x\mapsto\mathrm{rep}_{D,\cV}(x_{(v)})} of \m{\GL_{I(D)}(\F_v)}: $L$ is
injective by (i), \m{e_{I(D),c_j}} does not kill $y_j$, the parahoric type is (iii), the $\chi_n$ are distinct
by (iv), and no \m{x_{\chi_n}\in\F_v} is a root of $c_j'$, which is irreducible of degree \m{r>1}.

\emph{Step 3: $f_1$ and $f_2$ are eigenforms for one family of polynomials.} By Input~4(b) there are, for every
split prime $w$ with neither $w$ nor $\bar w$ chosen, numbers \m{t_{w,0},t_{w,1},\dots} that depend only on
$\Upsilon_w$ such that every function spherical at $w$ with $\Upsilon_w$ is an eigenvector of \m{\mathbb T^{w,n}}
with eigenvalue~$t_{w,n}$. Let \m{\mathrm{Hk}_w=\sum_n(-1)^n(\#\F_w)^{n(n-1)/2}\,t_{w,n}\,X^{\#I'-n}} for these
primes, and \m{\mathrm{Hk}_w=0}, say, for the others. By (ii), $f_1$ and $f_2$ have Hecke polynomial
$\mathrm{Hk}_w$ at every such~$w$. They are not zero, by the last condition of a product label. So both are
eigenforms for \m{(\mathrm{Hk}_w)_w}, by Lemma~\ref{l:eigen}.

\emph{Step 4: one Galois representation, two labellings.} By Input~1, applied to~$f_1$, there is a homomorphism
$\rho$ attached to \m{(\mathrm{Hk}_w)_w}. Input~2 applies to $\rho$ with $f_1$ and also with $f_2$, since it
holds for every attached representation. For \m{j\in\{1,2\}} let
\[
Q_j=\prod_{x:\,c_j'(x)=0}\bigl(X-\iota T(x)\bigr)\cdot\prod_{n=1}^{n_1}\bigl(X-\iota T(x_{\chi_n})\bigr)^{t_n},
\]
the first product over the roots of $c_j'$ in~$\bar{\F}_v$. Then
\begin{equation}\label{e:kill}
Q_1(\rho(\sigma))=0,\qquad Q_2(\rho(\sigma))=0,\qquad
\det\rho(\sigma)=\prod_{x:\,c_1'(x)=0}\iota T(x)\cdot\prod_{n=1}^{n_1}\iota T(x_{\chi_n})^{t_n}.
\end{equation}

\emph{Step 5: if \m{c_1'\ne c_2'}, every eigenvalue of \m{\rho(\sigma)} is a root of unity of order dividing
$M_\omega(D)$.} Let $\lambda$ be an eigenvalue. By \eqref{e:kill} it is a root of $Q_1$ and of~$Q_2$. Suppose
\m{\lambda=\iota T(x)} with $x$ a root of~$c_1'$. As $\iota$ and $T$ are injective and $\lambda$ is a root
of~$Q_2$, either $x$ is a root of~$c_2'$, or \m{x=x_{\chi_n}} for some~$n$. The first is impossible: two
different monic irreducible polynomials have no common root. So in every case \m{\lambda=\iota T(x_{\chi_n})} for
some~$n$, and \m{\lambda^{M_\omega(D)}=1} by (vi), because \m{\chi\mapsto x_\chi} is a homomorphism.

\emph{Step 6: conclusion.} Suppose \m{c_1'\ne c_2'}. The determinant of \m{\rho(\sigma)} is the product of its
eigenvalues, so \m{(\det\rho(\sigma))^{M_\omega(D)}=1} by Step~5. In the formula of \eqref{e:kill} the second
product also has $M_\omega(D)$-th power~$1$. Hence the product $\Pi$ of the roots of $c_1'$ satisfies
\m{\iota T(\Pi)^{M_\omega(D)}=1}, and \m{\Pi^{M_\omega(D)}=1} because $\iota$ and $T$ are injective and
multiplicative. But by Step~1 and Input~5(b), $\Pi$ is a primitive root of unity of order \m{q-1}, and
\m{0<M_\omega(D)<q-1}: a contradiction. So \m{c_1'=c_2'}, and \m{c_1=c_2}, since \m{c\mapsto c^{[x_\psi]}} is a
bijection.
\end{proof}

\sect{Where the results are used}{s:use}
Theorem~1 is an input of the note \emph{Good forms and rows}; it gives there that the space of the rows and the
sources of the single lifts have rigid labels, which is the hypothesis of the comparison of the traces of Hecke
operators on good forms. Lemma~\ref{l:eigen} is used once more, in Chapter~VIII, where rows are shown to give
eigenforms.

\sect{Notes on the formal text}{s:notes}
\begin{enumerate}[label=\textup{\arabic*.},leftmargin=1.6em,itemsep=5pt]
\item \emph{Letters.} The formal text writes {\leanfont x} for the chosen prime~$v$, {\leanfont v} for the
injection~$\nu$, {\leanfont n} for~$I'$, {\leanfont m} and {\leanfont b} for $n_1$ and~$\beta$, {\leanfont F} for
the family~$\Upsilon$, {\leanfont Q} for the family of Hecke polynomials, {\leanfont M} for the set~$\mathcal W$,
{\leanfont r} for the Galois representation~$\rho$, {\leanfont s} for the inertia element~$\sigma$, {\leanfont A}
for the place~$\mathfrak P$, and {\leanfont Ω} for~$\bar{\mathbb Q}$.
\item \emph{The functions $\Upsilon_w$.} In (ii) the value $\Upsilon_w$ is a function of a complex variable, not a
list of eigenvalues: the hypothesis prescribes, for every complex~$X$, the scalar by which the Hecke polynomial
evaluated at $X$ acts. That the family is common to all forms of the eigenpart is what makes Step~3 work.
\item \emph{The target of $L$.} The weight, the system of places, the subgroup and the character that determine
the space of algebraic forms are implicit arguments; the proof uses only that its elements are algebraic modular
forms on the standard space, so that Inputs 1 and~2 apply.
\item \emph{The prime $2$.} The formal proof fixes \m{\ell=2} and obtains the isomorphism $\iota$ from a lemma of
the library on algebraically closed fields of the same cardinality (the catalogue names Steinitz; the second
reading derives it from the existence of transcendence bases and the uniqueness of algebraic closures). Any
prime other than the residue characteristic of $v$ would do.
\item \emph{The end of the proof.} The formal text closes Steps 5 and~6 with two automatic tactics after a
rewriting that turns the claim \m{\Pi^{M_\omega(D)}=1} into the statement that every root of the characteristic
polynomial of \m{\rho(\sigma)} has $M_\omega(D)$-th power~$1$; the roots of the characteristic polynomial enter as
the multiset whose product is the determinant, and each of them is shown to be a root of $Q_1$ and of~$Q_2$. I
read these lines as Steps 5 and~6 and did not replay the goals. The quantified objects are the eigenvalues of
\m{\rho(\sigma)}: the roots of its characteristic polynomial, in lines 51 and 58 to~59 of the unit's file. The roots
of $c_1'$ and~$c_2'$ enter twice: as the case excluded by irreducibility in Step~5, and through their product, by
the determinant, in Step~6. An independent referee of this chapter of the handout read the same lines term by term and confirms this
reading, and the dossier of the unit, in its corrected state, follows it. The second reader describes the
contradiction as a root of~$c_1'$ falling among the roots \m{x_{\chi_n}}; read literally, that route would need
every root of~$Q_1$ to be an eigenvalue of \m{\rho(\sigma)}, which Input~2 does not give. The determinant clause of
Input~2 is essential: that $Q_1$ and $Q_2$ kill \m{\rho(\sigma)} would not suffice.
\item \emph{Lemma~\ref{l:eigen}.} In the formal statement the set $\mathcal W$ is an arbitrary type and $I'$ need
not be non-empty. The family of polynomials is only constrained at the split primes off the level.
\item \emph{What is not used.} The hypotheses (v) and (vi) enter only through the orders of $x_\psi$ and
of the~$x_{\chi_n}$. The embedding $\nu$ and the map $\beta$ enter only through Input~3.
\end{enumerate}

\subhead{How this note was written.} From the dossier of the unit: its two statements and its two structured
proofs, which a program compared line by line with the source (\UnitSteps{} leaf steps), and the second reading
of the unit by the other agent of this write-up. I read the Lean file of the unit completely, in particular its
last thirteen lines term by term, and in the library the definitions of Hecke polynomials, eigenforms, spherical
vectors, parahoric types, product labels, tame exponents, attached representations and the predicate that a
polynomial kills an element, and the statements of Inputs 1 to~4. The two statements of this note were compared
with their Lean statements argument by argument, and the order of the quantifiers is that of the formal text.
The proof follows the formal one; its division into six steps is mine. The references are those of the
library's files, of the catalogue of this handout and of its second reading; I did not open the cited works.

\head{Appendix A. Concordance with the Lean artifact}
{\fontsize{8.6}{11.6}\selectfont\renewcommand{\leanfont}{\ttfamily\fontsize{7.7}{10}\selectfont\color{leancol}}
\setlength{\LTleft}{0pt}\setlength{\LTright}{\fill}\setlength{\tabcolsep}{0pt}
\begin{longtable}{@{}p{30mm}p{56mm}p{43mm}r@{}}
\textbf{In this note} & \textbf{Lean declaration} & \textbf{File} & \textbf{Lines}\\\midrule\endhead
\ConcordanceRows
\end{longtable}}

\subhead{The two statements, and the two deep inputs, as they are written in the artifact.}
\quoted{isEigenformofhasHeckePolyAt}
\quoted{occurringLabelssubsingleton}
\quoted{sphericalCompatstd}
\quoted{isLabelledAt}
The first two stand after a line that declares the variables {\leanfont D}, {\leanfont s}, {\leanfont V} and the
finite type {\leanfont n}. In these texts {\leanfont IsSphericalWith}, {\leanfont HasParahoricType},
{\leanfont HasProductLabel}, {\leanfont IsEigenform}, {\leanfont Attached} are the notions (ii), (iii), (f), (c)
and (h) above; {\leanfont IsLabelledAt} with {\leanfont PolyKillsAt} is the conclusion of Input~2;
{\leanfont SphericalCompat}, in its first half, is Input~1; {\leanfont P.Mω D.s} is $M_\omega(D)$.

\head{Notation}
\begingroup\fontsize{8.6}{10.9}\selectfont\raggedright
\setlength{\columnsep}{9mm}\setlength{\parskip}{0pt}
\begin{multicols}{2}
\nt{$P=(N_0,d)$, $D$}{the parameters; admissible data (Setting~\ref{s:data})}
\nt{$E$, $E(D)$, $\bar{\mathbb Q}$}{a CM field; the field of the data; an algebraic closure of $\mathbb Q$}
\nt{$w$, $v$, $\bar w$, $\F_w$, $E_w$, $O_w$}{primes ($v$ a chosen one), the conjugate prime, residue field, completion, its integers}
\nt{$r(D)$, $m(D)$, $I(D)$, $\mathsf G(D)$, $M_\omega(D)$}{the rank, the number of lifts, the row index set, the finite group and the integer of the data}
\nt{$\fs$, $\cV$}{a family of signs; a space with a point}
\nt{$\mathrm{Prin}(D,\cV)$, $\mathrm{rep}_{D,\cV}$, $\mathrm{Prin}(D,\cV)[p]$, $\mathrm{Occ}(D,\cV,p,v)$}{principal forms, the action of $\mathsf G(D)$, eigenparts, occurring primitive labels (note \emph{Good forms and rows})}
\nt{$p$}{an eigensystem}
\nt{$I$, $I'$, $\GL_{I'}$, $G_f(I')$, $K_0(I')$, $g^*$}{finite index sets; matrix groups; the unitary group over the finite adeles and its integral points; conjugate transpose}
\nt{$k^{(w)}$, $\bar k$, $\bar k_{nn}$, $x_{(v)}$, $\tilde x$}{the element with component $k$ at $w$; reduction; its diagonal blocks; the element with component $x$ at $v$; a lift}
\nt{$R(g)$, $f$, $f_j$, $\mathcal W$}{right translation; functions on $G_f(I')$ with values in functions on $\mathcal W$}
\nt{$\mathbb T^{w,n}$, $\mathrm{Hk}_w$, $Q$, $t_n$, $t_{w,n}$}{Hecke operators, Hecke polynomials, eigenvalues}
\nt{$L$, $\nu$, $\beta$, $n_1$, $\Upsilon_w$, $\Upsilon$}{the data of Theorem~1}
\nt{$\psi$, $\chi$, $\chi_n$, $x_\chi$, $x_\psi$}{multiplicative characters of $\F_v$ and the units that name them}
\nt{$\F$, $\bar{\F}$, $T$, $\theta_x$, $\Xi_c$, $\Xi$, $e_{I,c}$, $c^{[u]}$}{a finite field, its closure, the injective character, the character named by~$x$, Green's class functions, idempotents, scaled polynomials (note \emph{Primitive labels})}
\nt{$c$, $c_j$, $c_j'$, $X$, $t=(t_n)$, $\delta$, $M$, $u$}{labels and their twists; the variable; sizes of blocks; a degree or level; an exponent; a unit}
\nt{$\kappa$, $\mathrm{alg}_\kappa$, $\gamma$}{a weight, its representation, a rational point}
\nt{$\ell$, $\iota$, $\bar{\mathbb Q}_\ell$}{a prime, an isomorphism \m{\mathbb C\to\bar{\mathbb Q}_\ell}, an algebraic closure of the $\ell$-adic numbers}
\nt{$\mathfrak P$, $\sigma$, $\varpi$, $z$, $\zeta$, $a$, $q$}{a place of $\bar{\mathbb Q}$, a Galois element, a uniformizer, a root of it, a root of unity; a tame exponent; the order of a finite field}
\nt{$\rho$, $Q_j$, $\lambda$, $\Pi$, $\eta$, $U$}{a Galois representation, the polynomials of the two labellings, an eigenvalue, the product of the roots of $c_1'$; a precision; an open set}
\nt{$y$, $y_j$, \m{x\cdot y}}{vectors, forms; the action of a finite group}
\nt{$n$, $j$, $i$, $x$, $g$, $k$, $b$, $\iota'$, $\pi$}{indices, elements of fields and groups; in \m{2\pi i} the circle constant and the imaginary unit}
\end{multicols}
\endgroup

\head{References}
{\fontsize{9.4}{12.6}\selectfont
All references are given as the library's files, the catalogue of this handout and its second reading record
them; I did not open the works.\par
T.~Barnet-Lamb, T.~Gee, D.~Geraghty, R.~Taylor, \emph{Potential automorphy and change of weight}, Ann. of Math.
179 (2014); Theorem~2.1.1.\par
T.~Barnet-Lamb, T.~Gee, D.~Geraghty, R.~Taylor, \emph{Local-global compatibility for $l=p$,~II}, Ann. Sci. ENS 47 (2014); Theorem~1.1 and
Corollary~1.2.\par
G.~Chenevier, M.~Harris, \emph{Construction of automorphic Galois representations,~II}, Camb. J. Math. 1 (2013);
Theorem~3.2.3.\par
M.~Harris, R.~Taylor, \emph{The geometry and cohomology of some simple Shimura varieties}, 2001.\par
J.-P.~Labesse, \emph{Changement de base CM et séries discrètes}, 2011; Corollaire~5.3.\par
J.~S.~Milne, \emph{Class Field Theory}, Chapter~VIII, Theorem~7.4, for the theorem of Chebotarev.\par
R.~Taylor, T.~Yoshida, \emph{Compatibility of local and global Langlands correspondences}, J. Amer. Math. Soc. 20
(2007); Theorem~A.\par}

\end{document}
